\documentclass{article}
\usepackage{graphicx} % Required for inserting images


% -------------------- Packages --------------------
\usepackage[utf8]{inputenc}
\usepackage[T1]{fontenc}
\usepackage{lmodern}          % Better font
\usepackage{geometry}         % Margins
\geometry{margin=1in}

\usepackage{amsmath, amssymb, amsthm}   % Math packages + theorem env
\usepackage{mathtools}        % Enhancements for amsmath
\usepackage{enumitem}         % Better lists
\usepackage{graphicx}         % Images
\usepackage{booktabs}         % Good tables
\usepackage{hyperref}         % Clickable refs
\usepackage{xcolor}           % Colors, if needed
\usepackage{tikz}             % Graph drawing (optional)
\usetikzlibrary{graphs}       % Graph theory tools (optional)

% -------------------- Theorem Environments --------------------
\theoremstyle{plain}
\newtheorem{theorem}{Theorem}
\newtheorem{lemma}{Lemma}
\newtheorem{proposition}{Proposition}
\newtheorem{corollary}{Corollary}

\theoremstyle{definition}
\newtheorem{definition}{Definition}
\newtheorem{example}{Example}

\theoremstyle{remark}
\newtheorem{remark}{Remark}

% Proof environment already included in amsthm.

% -------------------- Custom Commands --------------------
\newcommand{\N}{\mathbb{N}}
\newcommand{\Z}{\mathbb{Z}}
\newcommand{\Q}{\mathbb{Q}}
\newcommand{\R}{\mathbb{R}}
\newcommand{\C}{\mathbb{C}}

% Handy graph notation
\newcommand{\vn}{V(G)}
\newcommand{\en}{E(G)}

% ---------------------------------------------------------

\title{}
\author{}
\date{}

\begin{document}

\maketitle

\section{Functions and Relations}
\begin{proposition}
If $n$ is a positive integer, then $133$ divides $11^{n+1}+12^{2n-1}$.
\end{proposition}

\begin{proof}
Note that $133=7\cdot 19$ and $\gcd(7,19)=1$. It suffices to show
\[
11^{n+1}+12^{2n-1}\equiv 0 \pmod 7
\quad\text{and}\quad
11^{n+1}+12^{2n-1}\equiv 0 \pmod{19}.
\]

\textbf{Modulo $7$.}
We have $11\equiv 4\pmod 7$ and $12\equiv 5\pmod 7$. Thus
\[
11^{n+1}+12^{2n-1}
\equiv 4^{\,n+1} + 5^{\,2n-1} \pmod 7.
\]
Since $5\equiv -2$ and $25\equiv 4\pmod 7$,
\[
5^{2n-1} = 5\,(5^2)^{n-1} \equiv 5\cdot 4^{\,n-1}\pmod 7.
\]
Hence
\[
4^{\,n+1}+5^{\,2n-1}
\equiv 4^{\,n+1}+5\cdot 4^{\,n-1}
=4^{\,n-1}(4^2+5)
=4^{\,n-1}(16+5)
=4^{\,n-1}\cdot 21 \equiv 0 \pmod 7.
\]

\textbf{Modulo $19$.}
We have $12^2=144\equiv 11\pmod{19}$, so $11\equiv 12^2\pmod{19}$ and hence
\[
11^{n+1} \equiv (12^2)^{n+1} = 12^{2n+2}\pmod{19}.
\]
Therefore
\[
11^{n+1}+12^{2n-1}
\equiv 12^{2n+2}+12^{2n-1}
=12^{2n-1}\bigl(12^3+1\bigr)\pmod{19}.
\]
Compute $12^3\pmod{19}$:
\[
12^2\equiv 11,\quad 12^3\equiv 11\cdot 12 = 132 \equiv -1 \pmod{19}.
\]
Thus $12^3+1\equiv -1+1\equiv 0\pmod{19}$, and so
\[
11^{n+1}+12^{2n-1}\equiv 12^{2n-1}\cdot 0 \equiv 0\pmod{19}.
\]

We have shown that $11^{n+1}+12^{2n-1}$ is divisible by both $7$ and $19$.
Since $\gcd(7,19)=1$, it follows that $11^{n+1}+12^{2n-1}$ is divisible by
$7\cdot 19=133$. Hence $133\mid 11^{n+1}+12^{2n-1}$ for every positive
integer $n$.
\end{proof}

\begin{proposition}
    Define \(R\subseteq\mathbb{R}\times\mathbb{R}\) by
\((a,b)\in R \iff a-b\in\mathbb{Q}\).
\end{proposition}
\begin{proof}
    


\textbf{(i) $R$ is an equivalence relation.} \\
Reflexive: \(a-a=0\in\mathbb{Q}\Rightarrow (a,a)\in R\).\\
Symmetric: \(a-b\in\mathbb{Q}\Rightarrow b-a=-(a-b)\in\mathbb{Q}\Rightarrow (b,a)\in R\).\\
Transitive: \(a-b\in\mathbb{Q},\; b-c\in\mathbb{Q}\Rightarrow a-c=(a-b)+(b-c)\in\mathbb{Q}\Rightarrow (a,c)\in R\).

Hence \(R\) is an equivalence relation.

\medskip
\textbf{(ii) Cardinality of equivalence classes.} \\
The class of \(a\) is \( [a]=a+\mathbb{Q}=\{a+q:q\in\mathbb{Q}\}\). Since \(\mathbb{Q}\) is countable, each class \([a]\) is countable.

If the set of classes were countable, then \(\mathbb{R}\), being the union of all classes, would be a countable union of countable sets and therefore countable — contradiction. Thus the set of equivalence classes is uncountable.

\end{proof}

\begin{proposition}
There do not exist positive integers \(x,y\) satisfying
\[
x^{2}=2y^{2}+3.
\]
\end{proposition}

\begin{proof}
Reduce modulo \(4\). Squares are \(0\) or \(1\pmod4\). If \(y\) were even then
\(2y^{2}\equiv0\pmod4\) and hence \(x^{2}=2y^{2}+3\equiv3\pmod4\), impossible.
Thus \(y\) is odd. Write \(y=2k+1\). Then
\[
y^{2}=(2k+1)^{2}=4k(k+1)+1\equiv1\pmod8
\]
because \(k(k+1)\) is even, so \(4k(k+1)\) is a multiple of \(8\). Hence
\[
2y^{2}+3\equiv 2\cdot1+3=5\pmod8.
\]
But no square is congruent to \(5\pmod8\) (odd squares are \(1\pmod8\), even
squares are \(0\) or \(4\pmod8\)). This contradiction shows there are no
positive integers \(x,y\) with \(x^{2}=2y^{2}+3\).
\end{proof}
\begin{proposition}
There are no positive integers \(x,y\) satisfying
\[
x^2 = 2y^2 + 3.
\]
\end{proposition}

\begin{proof}
Work modulo \(3\). Squares mod \(3\) are \(0\) or \(1\).

Reduce the equation mod \(3\). Since \(3\equiv0\pmod3\) we obtain
\[
x^2 \equiv 2y^2 \pmod3.
\]

Case 1: \(y\not\equiv 0\pmod3\). Then \(y^2\equiv1\pmod3\), so
\(x^2\equiv 2\pmod3\), impossible because a square modulo \(3\) cannot be \(2\).

Case 2: \(y\equiv 0\pmod3\). Write \(y=3k\). Then
\[
x^2 = 2(3k)^2+3 = 18k^2+3 = 3(6k^2+1),
\]
so \(3\mid x^2\) and hence \(3\mid x\). Put \(x=3t\). Substituting and
dividing by \(3\) yields
\[
3t^2 = 6k^2+1 \quad\Longrightarrow\quad 0\equiv 1\pmod3,
\]
a contradiction.

Both cases are impossible. Therefore no positive integers \(x,y\) satisfy
\(x^2=2y^2+3\).
\end{proof}
\begin{lemma}
For any integer $n$, the residue of $n^2 \pmod 3$ is either $0$ or $1$.
\end{lemma}

\begin{proof}
Every integer $n$ is congruent modulo $3$ to exactly one of
\[
0,\;1,\;2.
\]
Compute the squares of these representatives:
\[
0^2 \equiv 0 \pmod 3,\qquad
1^2 \equiv 1 \pmod 3,\qquad
2^2 = 4 \equiv 1 \pmod 3.
\]
Thus every square is congruent to either $0$ or $1$ modulo $3$.
\end{proof}
\begin{proposition}
For an integer $n$, the following are equivalent:

(i) $n$ is odd.

(ii) $n^2 \equiv 1 \pmod{8}$.

(iii) $n^4 \equiv 1 \pmod{16}$.
\end{proposition}

\begin{proof}
\textbf{(i $\Rightarrow$ ii).}
If $n$ is odd, write $n = 2k+1$. Then
\[
n^2 = (2k+1)^2 = 4k(k+1)+1.
\]
Since $k(k+1)$ is even, $4k(k+1)$ is a multiple of $8$, hence
\[
n^2 \equiv 1 \pmod{8}.
\]

\medskip
\textbf{(ii $\Rightarrow$ iii).}
Assume $n^2 \equiv 1 \pmod{8}$. Then $n^2 = 8m + 1$ for some $m\in\mathbb{Z}$.  
Square both sides:
\[
n^4 = (8m+1)^2 = 64m^2 + 16m + 1 \equiv 1 \pmod{16}.
\]

\medskip
\textbf{(iii $\Rightarrow$ i).}
Assume $n^4 \equiv 1 \pmod{16}$. Reduce modulo $2$:
\[
n^4 \equiv 1 \pmod 2.
\]

\textbf{(iii $\Rightarrow$ i).}
Assume $n^4\equiv 1\pmod{16}$. If $n$ were even, write $n=2k$; then
\[
n^4=(2k)^4=16k^4\equiv 0\pmod{16}.
\]
Hence an even $n$ cannot satisfy $n^4\equiv1\pmod{16}$. Therefore $n$ is odd.

\medskip
Thus (i), (ii), (iii) are equivalent.
\end{proof}
\begin{proposition}
For all real numbers \(x,y\) we have
\[
|x+y|\ge \bigl||x|-|y|\bigr|.
\]
\end{proposition}

\begin{proof}
We split into two cases.

\medskip
\noindent\textbf{Case 1:} \(|x|\ge |y|\). \\
Then \(|x|-|y|\ge0\). Using the triangle inequality in the form
\(|x|=|(x+y)-y|\le |x+y|+|y|\) we obtain
\[
|x|-|y| \le |x+y|.
\]
Hence \(\bigl||x|-|y|\bigr|=|x|-|y|\le |x+y|\).

\medskip
\noindent\textbf{Case 2:} \(|x|<|y|\). \\
Then \(|y|-|x|>0\). Apply the same triangle-inequality trick with the roles
of \(x\) and \(y\) interchanged:
\(|y|=|(x+y)-x|\le |x+y|+|x|\), whence
\[
|y|-|x| \le |x+y|.
\]
Thus \(\bigl||x|-|y|\bigr|=|y|-|x|\le |x+y|\).

\medskip
In both cases \(\bigl||x|-|y|\bigr|\le |x+y|\). Therefore
\(|x+y|\ge \bigl||x|-|y|\bigr|\) for all real \(x,y\).
\end{proof}
\begin{proposition}
Let $\mathcal I=\{A\subseteq\mathbb N: A\text{ is infinite}\}$ and define
$R$ on $\mathcal I$ by $A\,R\,B\iff A\subseteq B$.
\begin{enumerate}
  \item $R$ is a partial order on $\mathcal I$.
  \item $R$ is not a total (linear) order.
  \item $\mathcal I$ has no minimum element, but it has a maximum element (namely $\mathbb N$).
\end{enumerate}
\end{proposition}

\begin{proof}
(1) Reflexive: $A\subseteq A$ for every $A\in\mathcal I$.  
Antisymmetric: $A\subseteq B$ and $B\subseteq A$ imply $A=B$.  
Transitive: $A\subseteq B$ and $B\subseteq C$ imply $A\subseteq C$.  
Hence $R$ is a partial order.

\medskip\noindent
(2) Not total: take the infinite sets of even and odd naturals,
$E=\{2,4,6,\dots\}$ and $O=\{1,3,5,\dots\}$. Then $E,O\in\mathcal I$ but
neither $E\subseteq O$ nor $O\subseteq E$, so $R$ is not linear.

\medskip\noindent
(3) Minimum: suppose $M\in\mathcal I$ were least, i.e.\ $M\subseteq A$ for every
infinite $A$. Choose any $m\in M$. The set $A=\mathbb N\setminus\{m\}$ is
infinite, but $m\notin A$, hence $M\not\subseteq A$, contradiction. Thus no
minimum exists.

Maximum: $\mathbb N\in\mathcal I$ and for every $A\in\mathcal I$ we have
$A\subseteq\mathbb N$, so $\mathbb N$ is the greatest element (maximum).
\end{proof}
\section{Group and Graph}
\begin{proposition}
Let $G$ be a cyclic group of order $24$. Then $G$ has exactly $2$ elements of order $4$.
\end{proposition}

\begin{proof}
Since $G$ is cyclic of order $24$, let $G=\langle g\rangle$. Every element of $G$
is of the form $g^k$ for $0\le k\le 23$, and its order is
\[
\operatorname{ord}(g^k)=\frac{24}{\gcd(24,k)}.
\]

We seek all $k$ such that $\operatorname{ord}(g^k)=4$. Thus:
\[
\frac{24}{\gcd(24,k)}=4
\quad\Longrightarrow\quad
\gcd(24,k)=6.
\]

Now list all $k$ with $0\le k<24$ and $\gcd(24,k)=6$. The multiples of $6$
below $24$ are $0,6,12,18$. Their gcds with $24$ are:
\[
\gcd(24,0)=24,\qquad
\gcd(24,6)=6,\qquad
\gcd(24,12)=12,\qquad
\gcd(24,18)=6.
\]
Thus exactly two values satisfy $\gcd(24,k)=6$, namely $k=6$ and $k=18$.

Therefore the only elements of order $4$ are $g^6$ and $g^{18}$, and hence
there are exactly $2$ elements of order $4$ in $G$.
\end{proof}
% Q3-A
\begin{proposition}
The statement
``Every connected bipartite graph has at least one vertex of even degree''
is false.
\end{proposition}

\begin{proof}
Consider the complete bipartite graph $K_{3,3}$. It is connected and
bipartite with bipartition $\{a_1,a_2,a_3\}\cup\{b_1,b_2,b_3\}$, and
$\deg(a_i)=\deg(b_j)=3$ for all $i,j$. Thus every vertex has odd degree and
there is no vertex of even degree. Hence the given statement is false.
\end{proof}
\begin{proposition}
Let $G$ be a connected graph in which every block (i.e.\ maximal $2$-connected
subgraph) is bipartite. Then $G$ is bipartite.
\end{proposition}

\begin{proof}
Recall two standard facts:

(1) A cycle is $2$-connected, hence is contained in some block of $G$.

(2) A graph is bipartite $\iff$ it contains no cycle of odd length.

\medskip
Assume, for contradiction, that $G$ is not bipartite. By (2) $G$ then
contains an odd cycle $C$. By (1) the cycle $C$ is contained in some block
$B$ of $G$. Thus $B$ contains an odd cycle, so $B$ is not bipartite, which
contradicts the hypothesis that every block of $G$ is bipartite.

Hence $G$ has no odd cycle. Again by (2), $G$ is bipartite.
\end{proof}

\end{document}